\newpage
\thispagestyle{empty}
\addcontentsline{toc}{chapter}{Abstract}
\chapter*{\centerline{ABSTRACT}\vspace{-1cm}} % Adjust the value as needed

    Data Visualization is indispensable across various domains such as business, and healthcare as it enables decision-makers to quickly grasp the essence of data, facilitating faster and more informed decision-making processes. Data Visualization is the process of translating data into visual elements like charts, graphs, and maps. This process is done to represent the data in a way that is easier to process, understand, and interpret. Nowadays, data in business exploration and technological advancements are usually high-dimensional which may lead to many challenges. Some of these are 1) The Curse of Dimensionality 2) Issues with visualization and its interpretation 3) Preserving different information like Local structures, Global structures, and Clusters. These challenges can be addressed through Dimensionality Reduction (DR) techniques and the process ensures that the insights of data are preserved well.

Dimensionality Reduction (DR) techniques are of two categories, linear methods and nonlinear methods. Some of the Nonlinear Dimensionality Reduction (NLDR) methods are t-SNE \cite{Hinton2002}, UMAP \cite{umap_reference}, LE \cite{belkin2003laplacian}, Isomap \cite{tenenbaum2000global}, and KPCA \cite{scholkopf1998nonlinear}. In data visualization preserving potential clusters information is one of the challenges that can be addressed through the NLDR techniques. Clusters are important in various domains such as marketing, and medical fields enabling targeted strategies, risk management, and insights extraction. Clusters' significance lies in revealing patterns, enhancing decision-making, and customizing approaches for distinct groups within intricate datasets. Sun et al. \cite{sun2023laplacian} discovered an NLDR algorithm that is LaptSNE, which explores and exploits the potential clusters of the data and produces cluster-informative low-dimensional visualization. McInnes et al.  \cite{umap_reference} claimed that the UMAP has no computational restrictions, and it can scale to significantly larger data set sizes compared to the practical limits of tSNE. Authors claimed that the UMAP algorithm is a strong contender to tSNE to preserve more of the global structures and for visualization quality with superior run time performance. 

The works of Sun et al. \cite{sun2023laplacian} and McInnes et al. \cite{umap_reference} motivated us to propose LAPUMAP to improve visualization performance. The proposed NLDR approach (LAPUMAP) is based on UMAP and may preserve more of the global structures and generate high-quality embedding. The experiments are to be performed on real-world datasets.


